\section{Dyadic vector masses and common prices}
\label{sec:prices}

Throughout this section, let $g\in C_c^\infty(\R^2)$ be real valued, and write
\[
 h=\nabla g\,\dd x,
 \qquad V=|h|(\R^2),
 \qquad u_t=A_tg.
\]
The averages $u_t$ are signed. Fix an integer $m\ge1$; all estimates below
are uniform in $m$.
Lemma~\ref{lem:local} discounts a cancelling error according to the
size of its support. We now decompose the gradient into errors to which
that discount applies, and assign one summable price to each dyadic
cell. The same price will work for every direction in a fixed finite
palette. The construction starts with the actual density of each scalar
coordinate of the gradient at the finest level and removes equal positive
and negative masses as it moves toward coarser levels.

\subsection{Positive refinement and cone coordinates}

For every integer $k$, set $r_k=2^{-k}$ and $V_k=r_k\mathbb Z^2$.
These definitions are not restricted to $1,\ldots,m$: the neighboring
direction estimate in Section~\ref{sec:labels} will also use level
$k-12$.
A vertex includes its level, so vertices at different levels are distinct
even when their spatial positions agree. For $p\in V_k$, define
\[
 Q_p=p+[0,r_k)^2,
 \qquad
 \phi_p(x)=\prod_{j=1}^2
       \left(1-\frac{|x_j-p_j|}{r_k}\right)_+.
\]
For $p\in V_i$ and $q\in V_k$ with $i<k$, put
$\lambda_{pq}=\phi_p(q)$, and at equal levels put
$\lambda_{pq}=\mathbf 1_{\{p=q\}}$. These nonnegative coefficients satisfy
\begin{equation}
\begin{gathered}
 \sum_{p\in V_i}\lambda_{pq}=1,
 \qquad
 \phi_p=\sum_{q\in V_k}\lambda_{pq}\phi_q
 \quad (i\le k),
 \\
 \lambda_{ps}=\sum_{q\in V_j}\lambda_{pq}\lambda_{qs}
 \quad (i\le j\le k,\ p\in V_i,\ s\in V_k).
\end{gathered}
\label{eq:hat-refinement}
\end{equation}
To verify these identities, first consider one coordinate. The hats on a
fixed grid sum to one: between two consecutive vertices, the only nonzero
hats have complementary affine values. A coarse hat is affine on every
interval of any finer dyadic grid, since all its breakpoints belong to
that grid. It therefore equals the piecewise affine interpolant of its
values at the finer vertices. Taking products proves the first two
identities in \eqref{eq:hat-refinement}; evaluating the second identity at
$s$ proves the third. The sums are locally finite. Notice that the normalization is in
the coarse index $p$ for each fixed fine vertex $q$. These coefficients
are refinement weights, not stochastic transition rows.

If $q\in V_{k+1}$ and $\lambda_{pq}>0$, then in each coordinate
$q_j-p_j\in\{-r_k/2,0,r_k/2\}$. Consequently
\begin{equation}
 \supp\phi_q\subseteq\supp\phi_p.
 \label{eq:prices-support-nesting}
\end{equation}
Thus each vertex has nine positive-weight children, and the supports
remain nested along every path of positive-weight edges.

Define the vector masses and their polar decompositions by
\[
 h_p=\int\phi_p\,\dd h,
 \qquad a_p=|h_p|,
 \qquad h_p=a_p\theta_p,
\]
where $\theta_p$ is an arbitrary fixed unit vector when $a_p=0$.
Refinement gives
\[
 h_p=\sum_{q\in V_k}\lambda_{pq}h_q
 \quad(p\in V_i,\ i\le k).
\]
The triangle inequality and the column sums in
\eqref{eq:hat-refinement} imply
\begin{equation}
 \sum_{p\in V_i}a_p
 \le\sum_{q\in V_k}a_q
 \le V
 \quad(i\le k).
 \label{eq:prices-vector-mass-budget}
\end{equation}
The last inequality follows directly from $\sum_q\phi_q=1$.
At each fixed level only finitely many of these masses can be nonzero,
because $h$ has compact support.

Fix $\tau=1/100$ and a finite $\tau/20$-net $\mathcal N$ of the unit
circle, measured in chord distance. For $n\in\mathcal N$, let $n^\perp$
be its counterclockwise quarter turn and put
\[
 e_+(n)=\frac{\sqrt3}{2}n+\frac12n^\perp,
 \qquad
 e_-(n)=\frac{\sqrt3}{2}n-\frac12n^\perp.
\]
These are unit vectors and $e_+(n)\cdot e_-(n)=1/2$.
Call $(q,n)$ \emph{good} when $\theta_q$ belongs to the closed sector
between $e_-(n)$ and $e_+(n)$, and \emph{bad} otherwise.
The two scalar coefficient functionals
\begin{equation}
 c_{n,+}(z)=\frac{n\cdot z}{\sqrt3}+n^\perp\cdot z,
 \qquad
 c_{n,-}(z)=\frac{n\cdot z}{\sqrt3}-n^\perp\cdot z
 \label{eq:prices-scalar-coordinates}
\end{equation}
satisfy
\[
 z=e_+(n)c_{n,+}(z)+e_-(n)c_{n,-}(z),
 \qquad
 |c_{n,+}(z)|+|c_{n,-}(z)|
 =2\max\left\{\frac{|n\cdot z|}{\sqrt3},
                     |n^\perp\cdot z|\right\}
 \le2|z|.
\]
In particular, for a good pair $(q,n)$, both coefficients of $h_q$ are
nonnegative, including when $h_q=0$. This sign information will supply
the positive cone measures in Lemma~\ref{lem:local}.

For the finite-band assembly, we also need to bound a vector's norm
by its signed projection along $n$ and the unfavorable projections
controlled by the local estimate. The following inequality gives
that bound, with an absolute constant:
\begin{equation}
 |z|\le C\left(n\cdot z+
                \sum_{\varepsilon\in\{+,-\}}
                (-e_\varepsilon(n)\cdot z)_+\right)
 \quad(z\in\R^2).
 \label{eq:cone-algebra}
\end{equation}
Here the expression in parentheses is nonnegative. Indeed, writing
$a=n\cdot z$, its value is positive for every nonzero $z$: if $a>0$
this is immediate; if $a=0$, one of the two negative parts equals
$|z|/2$; and if $a<0$, their sum is at least
$-(e_+(n)+e_-(n))\cdot z=\sqrt3|a|$, so the expression is at least
$(\sqrt3-1)|a|>0$. Its minimum on the unit circle is therefore positive
by continuity. Rotation invariance makes this minimum independent of
$n$, and homogeneity proves \eqref{eq:cone-algebra}.

\subsection{Cancellation from the finest level upward}

Let $\mathcal C$ be the indexed finite family consisting of
$c_{n,+},c_{n,-}$ for every $n\in\mathcal N$; repeated functionals,
if any, retain their indices. For $c\in\mathcal C$ write
\[
 \nu^c=c(h)=c(\nabla g(x))\,\dd x,
 \qquad c_s=c(h_s)=\int\phi_s\,\dd\nu^c.
\]
For $s\in V_l$, $1\le l\le m$, define
\begin{equation}
 d_s^c=
 \begin{cases}
 \displaystyle\frac12\left(
       \sum_{t\in V_{l+1}}\lambda_{st}|c_t|-|c_s|\right),
       & l<m,\\[6pt]
 \displaystyle\frac12\left(
       \int\phi_s(x)|c(\nabla g(x))|\,\dd x-|c_s|\right),
       & l=m.
 \end{cases}
 \label{eq:prices-scalar-defects}
\end{equation}
These numbers are nonnegative by refinement and the triangle
inequality. They measure the equal positive and negative masses
that can be removed at $s$. At level $m$ the definition uses the
positive and negative parts of the density itself. In particular,
$c_s=0$ need not imply $d_s^c=0$.

We will realize this numerical cancellation by measures supported
in the corresponding tent. The elementary operation
is proportional removal; it does not require the two measures
being paired to be mutually singular.

\begin{lemma}[Removing equal masses]\label{lem:prices-remove}
Let $U,V$ be finite nonnegative measures, of masses $P,N$.
There are nonnegative submeasures $\mu_+\le U$, $\mu_-\le V$,
both of mass $d=\min(P,N)$, such that
\[
 U-V=\rho+\mu_+-\mu_-,\qquad
 \rho(\R^2)=P-N,\qquad |\rho|(\R^2)=|P-N|.
\]
The residual $\rho$ has constant sign. If $U,V$ are absolutely
continuous and supported in a common closed set, so are all the
constructed measures.
\end{lemma}
\begin{proof}
When $P>0$ set $\mu_+=(d/P)U$; when $P=0$ set $\mu_+=0$.
Define $\mu_-$ in the same way from $N,V$, and set
$\rho=(U-\mu_+)-(V-\mu_-)$. If $P\ge N$ and $P>0$, then
$\mu_-=V$ and $\rho=(1-N/P)U\ge0$. The case $N\ge P$ is
symmetric. If $P=N=0$, every measure is zero. These formulas
prove the mass, sign and support assertions, including all zero
cases.
\end{proof}

\begin{lemma}[Exact scalar pairing]\label{lem:prices-scalar-pairing}
For every $c\in\mathcal C$ there are measures
$\rho_s^c,\mu_{s,+}^c,\mu_{s,-}^c$ at all vertices
$s\in V_l$, $1\le l\le m$, supported in $\supp\phi_s$ and
absolutely continuous, with $\rho_s^c$ of constant sign and
$\mu_{s,\pm}^c$ nonnegative, such that
\begin{equation}\label{eq:prices-scalar-masses}
 \rho_s^c(\R^2)=c_s,\qquad |\rho_s^c|(\R^2)=|c_s|,
 \qquad
 \mu_{s,+}^c(\R^2)=\mu_{s,-}^c(\R^2)=d_s^c.
\end{equation}
Put $\beta_s^c=\mu_{s,+}^c-\mu_{s,-}^c$. For every
$q\in V_k$, $1\le k\le m$, the following is an identity of
measures on the whole plane:
\begin{equation}\label{eq:prices-vertex-reconstruction}
 \phi_q\nu^c=\rho_q^c+
       \sum_{l=k}^m\sum_{s\in V_l}\lambda_{qs}\beta_s^c.
\end{equation}
Only finitely many terms in this identity are nonzero, and
\begin{equation}\label{eq:prices-pair-properties}
 \beta_s^c(\R^2)=0,\qquad
 |\beta_s^c|(\R^2)\le2d_s^c,\qquad
 \supp\beta_s^c\subseteq\supp\phi_s
       \subseteq\overline B(s,\sqrt2 r_l).
\end{equation}
\end{lemma}

\begin{proof}
Fix $c$ and first take $s\in V_m$. Apply
Lemma~\ref{lem:prices-remove} to the actual terminal measures
\[
 U_s=\phi_s(\nu^c)_+,
 \qquad V_s=\phi_s(\nu^c)_-.
\]
Their masses have difference $c_s$ and sum
$\int\phi_s\,\dd|\nu^c|$. Hence their smaller mass is exactly
$d_s^c$, and removal gives \eqref{eq:prices-scalar-masses} and
\begin{equation}\label{eq:prices-terminal-identity}
 \phi_s\nu^c=\rho_s^c+\beta_s^c.
\end{equation}
If the positive and negative terminal masses agree, the residual
is zero while both paired measures keep that common mass.

Suppose that all residuals at level $l+1$ have been constructed.
At $s\in V_l$ combine their positive and negative parts separately:
\[
 U_s=\sum_{t\in V_{l+1}}\lambda_{st}(\rho_t^c)_+,
 \qquad
 V_s=\sum_{t\in V_{l+1}}\lambda_{st}(\rho_t^c)_-.
\]
The residuals have constant sign, so the masses $P_s,N_s$ of
these two measures obey
\[
 P_s-N_s=\sum_t\lambda_{st}c_t=c_s,
 \qquad
 P_s+N_s=\sum_t\lambda_{st}|c_t|.
\]
Thus $\min(P_s,N_s)=d_s^c$. Applying
Lemma~\ref{lem:prices-remove} again constructs the measures in
\eqref{eq:prices-scalar-masses} and proves
\begin{equation}\label{eq:prices-residual-recursion}
 \sum_{t\in V_{l+1}}\lambda_{st}\rho_t^c
      =\rho_s^c+\beta_s^c.
\end{equation}
Absolute continuity is preserved by finite sums and scalar
multiplication. Support nesting in
\eqref{eq:prices-support-nesting} keeps all these measures in
$\supp\phi_s$. The construction proceeds through finitely many
levels. At each level only finitely many measures can be nonzero,
since the terminal density has compact support and each vertex
has finitely many parents.

It remains to verify that no density has been lost during removal.
At level $m$, \eqref{eq:prices-vertex-reconstruction} is
\eqref{eq:prices-terminal-identity}. If the identity holds for
all children of $q\in V_k$, refinement yields
\begin{align*}
 \phi_q\nu^c
 &=\sum_{t\in V_{k+1}}\lambda_{qt}\phi_t\nu^c\\
 &=\sum_t\lambda_{qt}\rho_t^c
   +\sum_{l=k+1}^m\sum_{s\in V_l}
        \left(\sum_t\lambda_{qt}\lambda_{ts}\right)\beta_s^c\\
 &=\rho_q^c+\beta_q^c
   +\sum_{l=k+1}^m\sum_{s\in V_l}\lambda_{qs}\beta_s^c.
\end{align*}
The last line uses \eqref{eq:prices-residual-recursion} and
coefficient composition in \eqref{eq:hat-refinement}. It is the
required identity at $q$. All rearrangements are finite for a
fixed root because branching and depth are finite. Finally,
equal masses in \eqref{eq:prices-scalar-masses} give zero mass
and variation at most $2d_s^c$ for each difference; the tent
support gives the last assertion in
\eqref{eq:prices-pair-properties}.
\end{proof}

The measures are constructed once at each vertex. A descendant's
pair may appear in several coarse identities, with its specified
coefficient; these identities do not assert a disjoint allocation.

\paragraph{The equivalent path form.}
The price estimate below uses the vertex identity. The following
equivalent form expresses shared-descendant coefficients as sums of
path weights, using the same cancelling pieces. Section~\ref{sec:labels}
will construct a separate positive measure on vector paths. Fix $q\in V_k$ with $1\le k\le m$, and let $\xi$ range
over all positive-weight paths $\xi=(q_k,\ldots,q_l)$ with
$q_k=q$ and $k\le l\le m$, including the trivial path $(q)$. Set
\[
 w_\xi=\prod_{j=k}^{l-1}\lambda_{q_jq_{j+1}},
 \qquad \mu_{\xi,\pm}^c=w_\xi\mu_{q_l,\pm}^c,
 \qquad w_{(q)}=1.
\]
Composition gives $\sum_{\xi:q_l=s}w_\xi=\lambda_{qs}$, and
therefore
\begin{equation}\label{eq:prices-scalar-pairing}
 \phi_q c(h)=\rho_q^c+
       \sum_\xi(\mu_{\xi,+}^c-\mu_{\xi,-}^c).
\end{equation}
Each path pair has mass $w_\xi d_s^c$ in each sign and support in
$\supp\phi_s$, where $s=q_l$. This recovers the path identity
without a separate allocation along every path.

The cancellation amounts have an exact total budget. Summing
\eqref{eq:prices-scalar-defects} over a level and using the column
normalization in \eqref{eq:hat-refinement} cancels intermediate
absolute masses. The terminal partition of unity gives
\begin{equation}
 \sum_{l=1}^m\sum_{s\in V_l}d_s^c
 =\frac12\left(
       \int|c(\nabla g)|\,\dd x-
       \sum_{s\in V_1}|c_s|\right)
 \le\frac12\|c\|\,V.
 \label{eq:prices-scalar-defect-budget}
\end{equation}
Here $\|c\|$ denotes the Euclidean operator norm, equal to
$2/\sqrt3$ for every $c\in\mathcal C$ by
\eqref{eq:prices-scalar-coordinates}. When $m=1$, the same formula
is just the sum of the terminal defects. Zero masses introduce no
exception to either the reconstruction or this telescoping identity.

\subsection{One price at every observation scale}

For $Q=Q_p$ at level $k$, put
\[
 \mathcal R(Q)=\{q\in V_k:|q-p|_\infty\le20r_k\},
 \qquad |z|_\infty=\max\{|z_1|,|z_2|\}.
\]
A pair created at level $l\ge k$ has support of radius at most
$\sqrt2r_l$. At observation scale $r_k$, the local estimate
therefore charges it the relative factor
$(r_l/r_k)^\alpha=2^{-\alpha(l-k)}$. We sum these charges over
all the scalar coordinates so that the resulting price is
independent of the label chosen later.

The scaled radius sets from Lemma~\ref{lem:local} are
\begin{align*}
 \mathcal I_k(x)
   &=\{r\in[r_k,2r_k]:u_r(x)\ge u_t(x)
             \text{ for every }t>0\text{ with }|t-r|\le r_k/4\},\\
 \mathcal J_k(x)
   &=\{r\in[r_k,2r_k]:u_r(x)=\max_{r_k\le t\le2r_k}u_t(x)\}
       \cup\{r_k,2r_k\}.
\end{align*}
Also define
\[
 v_k(x)=\sup_{0<t,t'\le4r_k}|u_t(x)-u_{t'}(x)|.
\]
An empty supremum of nonnegative quantities is zero, as before.

\begin{proposition}[Common discounted prices]
\label{prop:prices}
For $p\in V_k$, $1\le k\le m$, define
\begin{equation}
 D_{Q_p}=\sum_{c\in\mathcal C}
       \sum_{q\in\mathcal R(Q_p)}
       \sum_{l=k}^m2^{-\alpha(l-k)}
       \sum_{s\in V_l}\lambda_{qs}d_s^c,
 \qquad \alpha=\frac1{100}.
 \label{eq:prices-definition}
\end{equation}
These nonnegative prices satisfy
\begin{equation}
 \sum_{k=1}^m\sum_{p\in V_k}D_{Q_p}
 \le\frac{41^2}{2(1-2^{-\alpha})}
           \sum_{c\in\mathcal C}\|c(\nabla g)\|_{L^1}
 \le C V.
 \label{eq:price-budget}
\end{equation}
For every such cell $Q$ and every $n\in\mathcal N$,
\begin{equation}
\begin{aligned}
 &\int_Q\sup_{t\in\mathcal I_k(x)}
       \sum_{\varepsilon\in\{+,-\}}
       (-e_\varepsilon(n)\cdot\nabla u_t(x))_+\,\dd x
   \le C\left(D_Q+
       \sum_{\substack{q\in\mathcal R(Q)\\(q,n)\text{ bad}}}a_q\right),
 \\
 &\int_Q\left[\sup_{t\in\mathcal J_k(x)}|\nabla u_t(x)|
                  +r_k^{-1}v_k(x)\right]\dd x
   \le C\left(D_Q+\sum_{q\in\mathcal R(Q)}a_q\right).
\end{aligned}
\label{eq:prices}
\end{equation}
All constants are absolute: the direction palette, support radius
factor and cone aperture are fixed, and there is no dependence
on $m$.
\end{proposition}

\begin{proof}
We first sum the prices. At a fixed level each vertex $q$ belongs
to exactly $41^2$ neighborhoods $\mathcal R(Q_p)$, since each of
the two integer coordinate differences ranges from $-20$ to $20$.
Nonnegativity permits reordering, and the column sums give
\begin{align*}
 \sum_{k=1}^m\sum_{p\in V_k}D_{Q_p}
 &=41^2\sum_{c\in\mathcal C}\sum_{l=1}^m\sum_{s\in V_l}d_s^c
       \sum_{k=1}^l2^{-\alpha(l-k)}
          \sum_{q\in V_k}\lambda_{qs}\\
 &\le\frac{41^2}{1-2^{-\alpha}}
       \sum_{c\in\mathcal C}\sum_{l=1}^m\sum_{s\in V_l}d_s^c.
\end{align*}
Applying the exact telescoping identity
\eqref{eq:prices-scalar-defect-budget} proves
\eqref{eq:price-budget}. In particular, one may take
$C=41^2|\mathcal C|/[\sqrt3(1-2^{-\alpha})]$ in its final
inequality.

We next construct the decomposition needed by the local estimates.
Fix $Q=Q_p$ at level $k$, a label $n$, and write
$\Omega_Q=\{y:\dist(y,Q)<10r_k\}$. If a level-$k$ tent
$\phi_q$ is positive at $y\in\Omega_Q$, then
$|q-y|_\infty<r_k$ and $|y-p|_\infty<11r_k$, so
$q\in\mathcal R(Q)$. Hence
\begin{equation}
 h=\sum_{q\in\mathcal R(Q)}\phi_qh
 \quad\text{on }\Omega_Q.
 \label{eq:prices-local-partition}
\end{equation}
Apply \eqref{eq:prices-vertex-reconstruction} with the two
coordinates $c_{n,\varepsilon}$ and multiply by
$e_\varepsilon(n)$. This gives the measure identity
\begin{equation}\label{eq:prices-vector-decomposition}
 h=\sum_{q\in\mathcal R(Q)}\sum_{\varepsilon\in\{+,-\}}
            e_\varepsilon(n)\rho_q^{c_{n,\varepsilon}}
   +\sum_{q\in\mathcal R(Q)}\sum_{\varepsilon\in\{+,-\}}
        \sum_{l=k}^m\sum_{s\in V_l}
            e_\varepsilon(n)\lambda_{qs}\beta_s^{c_{n,\varepsilon}}
 \quad\text{on }\Omega_Q.
\end{equation}
The components here are measures on the whole plane. We retain
them on their full supports; only their sum is being compared
with $h$ on $\Omega_Q$.

For a good root $(q,n)$, both $c_{n,\varepsilon}(h_q)$ are
nonnegative. The constant-sign and mass assertions in
\eqref{eq:prices-scalar-masses} therefore imply
$\rho_q^{c_{n,\varepsilon}}\ge0$, including the zero case.
These residuals are the positive measures $\sigma_i$ of
Lemma~\ref{lem:local}, in the unit directions $e_+(n),e_-(n)$.
Each of these directions has dot product at least $1/2$ with
either test direction $e_+(n)$ or $e_-(n)$.
For a bad root, put the two signed residual terms among the
unrestricted errors, with assigned relative size one. Their total
variation is at most
\[
 \sum_{\varepsilon\in\{+,-\}}
       |\rho_q^{c_{n,\varepsilon}}|(\R^2)
 =\sum_{\varepsilon\in\{+,-\}}|c_{n,\varepsilon}(h_q)|
 \le2a_q.
\]
Thus no good-root mass enters the error price for the directional
estimate.

For each term in the second sum of
\eqref{eq:prices-vector-decomposition}, its total vector mass is
zero, its variation is at most
$2\lambda_{qs}d_s^{c_{n,\varepsilon}}$, and its support lies in
$\overline B(s,\sqrt2r_l)$. Assign it the relative size
$s_d=r_l/r_k=2^{k-l}$. When $l>k$ this gives exactly the
cancellation and support conditions required for a small error.
When $l=k$ it gives an allowed size-one error. The physical
sum of discounted variations is consequently bounded by
\begin{equation}
 \sum_d s_d^\alpha|\beta_d|(\R^2)
 \le2D_Q+
       2\sum_{\substack{q\in\mathcal R(Q)\\(q,n)\text{ bad}}}a_q.
 \label{eq:prices-directional-decomposition-cost}
\end{equation}
Here $\beta_d$ denotes the vector error pieces just specified.
The sum is finite and every component is absolutely continuous.
The terminal pairs are included, even at roots with zero vector
mass; there is no terminal remainder.

To apply Lemma~\ref{lem:local}, let $r=r_k$,
$T(X)=p+rX$ and $\widetilde g=g\circ T$. Its gradient measure
and disk averages satisfy
\begin{equation}\label{eq:prices-scaling}
 \widetilde h=r^{-1}(T^{-1})_\#h,
 \qquad A_t\widetilde g(X)=u_{rt}(T(X)),
 \qquad \nabla_XA_t\widetilde g(X)=r\nabla u_{rt}(T(X)).
\end{equation}
Indeed $\nabla_X\widetilde g=r\nabla g\circ T$ and
$\dd X=r^{-2}\dd x$. Transform every component measure by
the same $r^{-1}(T^{-1})_\#$ rule. Its mass and variation
are multiplied by $r^{-1}$; zero total vector mass and positivity
are preserved. A pair at level $l$ is now supported in a ball
of radius $\sqrt2r_l/r_k$. Taking $C_s=2$ places this closed
support inside an open ball of radius $C_s r_l/r_k$. The
transformed identity holds on the
$10$-neighborhood of $[0,1)^2$.

The radius sets for $\widetilde g$ are the images of
$\mathcal I_k,\mathcal J_k$ under division by $r$, including
the comparison with all positive radii in $\mathcal I_k$ and
both forced endpoints in $\mathcal J_k$. Furthermore, the
integral over the unit cell of any of the gradient suprema
in Lemma~\ref{lem:local} is $r^{-1}$ times the corresponding
physical integral over $Q$. Apply
\eqref{eq:local-direction} with cone margin $\kappa=1/2$
to each of the two test directions, and multiply by $r$.
The supremum of the sum of the two nonnegative terms is at
most the sum of their suprema. Together with
\eqref{eq:prices-directional-decomposition-cost}, this proves
the first estimate in \eqref{eq:prices}.

For the second estimate in \eqref{eq:prices}, fix any one label
and use the same decomposition, now placing \emph{all} root
residuals among the size-one errors. The positive part $S$
is zero, and the physical discounted error mass is at most
$2D_Q+2\sum_{q\in\mathcal R(Q)}a_q$. The size estimate
\eqref{eq:local-size}, after the same transformation, bounds
the physical gradient integral by an absolute multiple of
this quantity. For the value estimate the change of variables
instead gives
\begin{equation}\label{eq:prices-value-scaling}
 \int_{[0,1)^2}\sup_{0<t,t'\le4}
       |A_t\widetilde g(X)-A_{t'}\widetilde g(X)|\,\dd X
   =r^{-2}\int_Qv_k(x)\,\dd x.
\end{equation}
Its right side in \eqref{eq:local-value} is bounded by
$Cr^{-1}(D_Q+\sum_q a_q)$. Multiplying by $r$ therefore
bounds $r_k^{-1}\int_Qv_k$, exactly as asserted. Adding this
to the size bound completes the proof.
\end{proof}

The prices now pay for every scalar cancellation and leave only
nearby vector masses $a_q$ in the local bounds. Controlling those
remaining masses requires a choice of directions across scales
and positions. The next section constructs a positive measure on
complete refinement paths, recording a birth level, so that the mass
of paths through a vertex $q$ whose birth has already occurred is
exactly $a_q$. The scalar measures above supply neither this path
measure nor the spatial estimates for the direction law; their role
is the reconstruction and discounted budget proved in this section.
